\documentclass[exercice]{thedocument}%
\startdatakeys
    \source{15GENMATAN1}
    \BO{2026}
    \cycle{4}
    \competencesDuSocle{RA,CO}
    \objectifsApprentissage{B3ca}
    \motsCles{Thalès sens direct}
\enddatakeys
\begin{document}%
    Pour filmer les étapes d’une course cycliste, les réalisateurs de télévision
    utilisent des caméras installées sur deux motos et d’autres dans deux
    hélicoptères. Un avion relais, plus haut dans le ciel, recueille les images
    et joue le rôle d’une antenne relais. On considère que les deux hélicoptères
    se situent à la même altitude et que le peloton des coureurs roule sur une
    route horizontale. Le schéma ci-dessous illustre cette situation~:
    \begin{center}
        Insérer schéma
        https://www.planete-maths.fr/theoremethalesexercices3sujet.html
    \end{center}
    L’avion relais (point \mathspoint{A}), le premier hélicoptère (point
    \mathspoint{L}) et la première moto (point \mathspoint{N}) sont
    alignés.\medskip\par
    De la même manière, l’avion relais (point \mathspoint{A}), le deuxième
    hélicoptère (point \mathspoint{H}) et la deuxième moto (point \mathspoint{M})
    sont également alignés.\medskip\par
    On sait~: $\mathspoint{AM}=\mathspoint{AN}=1$ km ; $\mathspoint{HL}=270$ m et
    $\mathspoint{AH}=\mathspoint{AL}=720$ m.
    \begin{enumerate}
        \item Relever la phrase de l’énoncé qui permet d’affirmer que les droites
        \mathspoint{(LH)} et \mathspoint{(MN)} sont parallèles.
        \item Calculer la distance \mathspoint{MN} entre les deux motos.
    \end{enumerate}
\end{document}